\section{Introducci\`on}
El presente trabajo es una investigación y aplicaron de métodos preliminar para el manejo
de datos no lineales y su posterior análisis. El fin es entender como diferentes métodos no lineales
pueden ayudar a entender mejor los datos que estamos visualizando y de igual manera como podrían
ayudarnos a detectar tendencias o diferencias significativas entre nuestra población control y la
población de tri-atletas. Posteriormente se tiene previsto usar toda esta información y con ayuda de
los fractales que se han generado por medio de software observar las posibles tendencias o diferencias
mencionadas previamente, esto con el fin de justificar un análisis mas «reducido» (refiriéndonos a
comprobar la no linealidad y posterior análisis del fractal sin necesidad de pasar por demasiadas
pruebas de análisis o estadística) y así reducir considerablemente los tiempos de descripción de
fenómenos(en este caso enfermedades o identificación de sujetos)
